\documentclass[10pt,a4paper]{amsart}
\usepackage[margin=19mm]{geometry}
\usepackage{amsmath,amssymb,mathtools}
\usepackage[scr=rsfs,cal=euler]{mathalfa}
\usepackage{xfrac}
\usepackage[unicode,hidelinks,bookmarks=false]{hyperref}
\newcommand{\ie}{i.e.\ }
\newcommand{\cf}{cf.\ }
\newcommand{\resp}{respectively\ }
\newcommand{\Subsection}{\subsection}
\providecommand{\nobreakdash}{\nobreak\hbox{-}}
\setlength{\emergencystretch}{2em}
\multlinegap0pt
\usepackage{amssymb}
\newtheorem{proposition}{Proposition}
\newtheorem{lemma}{Lemma}
\newtheorem{corollary}{Corollary}
\newtheorem{claim}{Claim}
\title{Square-bracket operations clubs}
\author{Osvaldo Guzman, Stevo Todorcevic}
\date{Source: arXiv:2602.14373v1 (CC BY 4.0)}
\begin{document}
\maketitle

\noindent\emph{Source and licence.} Square-bracket operations clubs. Version arXiv:2602.14373v1 (CC BY 4.0), distributed by arXiv under the Creative Commons Attribution 4.0 International licence, http://creativecommons.org/licenses/by/4.0/, as recorded in the arXiv licence metadata for this exact version and shown on the abstract page https://arxiv.org/abs/2602.14373v1. Full licence notice: \texttt{licenses/arxiv-2602.14373-CC-BY-4.0.txt}.\par\smallskip

\noindent\textit{Editorial scope.} Counted unit: the article's proposition proposition (source0.tex lines 3732--3736). The article's complete original proof of it is delivered with it (source0.tex lines 3738--4187, one \texttt{proof} environment of 450 lines). No mathematics is written, repaired or re-derived: the statement and the proof are the article's own lines, in the article's own notation, and no retained line is substituted or re-flowed.  Retained uncounted as the source-local context the proof needs: lines 712--727; lines 745--749; lines 3570--3588.  Nothing was omitted from any retained span, so the package carries no omission record.  No retained line contains a citation, so the wrapper carries no bibliography and no bibliography entry.  No retained line contains a figure environment, an image-inclusion command or drawing code, so no image is delivered and the package declares no asset dependency.  Editorial formatting only, in addition: an AMS \texttt{amsart} preamble that restates the article's own declarations and macros used by the retained lines, namely the environments \texttt{proposition}, \texttt{lemma}, \texttt{corollary}, \texttt{claim} on separate counters, together with the standard packages those lines need.  The retained environments print in this document as Proposition 1, Lemma 1, Corollary 1, Claim 1 in order of appearance, whereas the article prints the same items with the numbering of its own full document.  The complete native source of the article as distributed in the arXiv e-print for this exact version is bundled unchanged as \texttt{sources/194706/source0.tex}.
\begin{proposition}
Let $\alpha<\gamma<\beta$ be countable ordinals. The following conditions are
equivalent: \label{Prop juntar trazas}

\begin{enumerate}
\item $\gamma\in$ \textsf{Tr}$\left(  \alpha,\beta\right)  .$

\item \textsf{Tr}$\left(  \alpha,\beta\right)  =$\textsf{Tr}$\left(
\gamma,\beta\right)  \cup$\textsf{Tr}$\left(  \alpha,\gamma\right)  .$

\item $\lambda\left(  \gamma,\beta\right)  <\alpha.$

\item $\rho_{0}\left(  \alpha,\beta\right)  =\rho_{0}\left(  \gamma
,\beta\right)  ^{\frown}\rho_{0}\left(  \alpha,\gamma\right)  .$
\end{enumerate}
\end{proposition}

\begin{lemma}
The set $\{\rho_{0\alpha}\mid\alpha\in$ \textsf{LIM}$\left(  \omega
_{1}\right)  \}$ is an antichain of $T\left(  \rho_{0}\right)  .$
\label{Lema limite anticadena}
\end{lemma}

\begin{corollary}
Let $T$ be an Aronszajn tree, $A\subseteq T^{n+1}$ and $\left\{
M_{0},...,M_{n}\right\}  $ an $\in$-chain of countable elementary submodels
such that $T,A\in M_{0}.$ Assume there is $\left(  s_{0},..,s_{n}\right)  \in
A$ such that $\delta_{M_{0}}\leq\left\vert s_{0}\right\vert <\delta_{M_{1}%
}\leq\left\vert s_{1}\right\vert \leq...\leq\delta_{M_{n}}\leq\left\vert
s_{n}\right\vert .$ There are $\gamma$ and $F$ with the following properties:
\label{Corol del juego}

\begin{enumerate}
\item $\gamma<\delta_{M_{0}}.$

\item $F=\left(  t_{0},...,t_{n}\right)  \in T^{n+1}$ and $t_{i}\neq
s_{j}\upharpoonright\gamma$ for every $i,j\leq n.$

\item Player $\mathsf{II}$ has a winning strategy in \emph{G}$_{F}\left(
T,A\right)  .$
\end{enumerate}
\end{corollary}

\begin{proposition}
Let $M$ be a big model and $p=(\mathcal{M}_{p},\mathcal{N}_{p})\in
\mathbb{P}\left(  C\right)  .$ If $\overline{M}\in\mathcal{N}_{p},$ then $p$
is $(M,\mathbb{P}\left(  C\right)  )$-generic.
\end{proposition}

\begin{proof}
For convenience, we will denote $T=T\left(  \rho_{0}\right)  .$ To simplify
notation, a subset $W\subseteq\omega_{1}$ will occasionally be identified with
the corresponding set $\{\rho_{0\alpha}\mid\alpha\in W\}.$

\qquad\ \qquad\ \qquad\ \qquad\ \qquad\ \qquad\ 

Let $D\in M$ be an open dense subset of $\mathbb{P}\left(  C\right)  $. By
extending $p$ if necessary, we may assume that $p\in D.$ We need to prove that
$p$ is compatible with an element of $D\cap M.$ Denote $\delta=\delta_{M},$
which is an element of $C.$ Define $X=$ \textsf{ht}$\left(  p\right)
\setminus M$ and let $m$ be the size of $X.$ Note that $\delta$ is the
smallest element of $X.$ Let $p_{M}=(\mathcal{M}_{p}\cap M,$ $\mathcal{N}%
_{p}\cap M).$ It is easy to see that $p_{M}\in\mathbb{P}\left(  C\right)  \cap
M$ and $p_{M}\sqsubseteq p.$

\qquad\ \qquad\ \qquad\ \ \ 

Find $\eta_{0}<\delta$ such that $S\left(  p\right)  \cap M\subseteq\eta_{0}.$
We now define $E$ as the set of all $Y\in\left[  \omega_{1}\right]  ^{m}$ for
which there is $r\in D$ with the following properties:

\begin{enumerate}
\item $p_{M}\sqsubseteq r.$

\item \textsf{ht}$\left(  r\right)  =$ \textsf{ht}$\left(  p_{M}\right)  \cup
Y.$

\item $S\left(  r\right)  \cap\eta_{0}=S\left(  p\right)  \cap\eta_{0}$ (which
is the same as $S\left(  p\right)  \cap M$).

\item $H\left(  p\right)  $ and $H\left(  r\right)  $ are isomorphic and if
$h:S\left(  p\right)  \longrightarrow S\left(  r\right)  $ denote the (unique)
isomorphism, then:

\begin{enumerate}
\item $h$ is the identity mapping on $S\left(  p\right)  \cap M.$

\item $\rho_{0\alpha}\upharpoonright\eta_{0}=\rho_{0h\left(  \alpha\right)
}\upharpoonright\eta_{0}$ for all $\alpha\in X.$
\end{enumerate}
\end{enumerate}

\qquad\ \qquad\ \qquad\ \ \ 

Note that $X\in E.$ Moreover, $E$ is in $M.$ This is because the definition of
$E$ can can be reformulated without reference to objects outside $M.$ For
$Y\in E$ and $r$ as above, we will say that $r$ is a witness of $Y\in E.$

\begin{claim}
There are $\eta_{1}$ and $F=\left(  t_{0},...,t_{m-1}\right)  $ with the
following properties:

\begin{enumerate}
\item $\eta_{1}\in M$ and $\eta_{0}<\eta_{1}<\delta.$

\item $F=\left(  t_{0},...,t_{m-1}\right)  \in\left(  T_{\eta_{1}}\right)
^{m}$ and $t_{i}\neq\rho_{0\alpha}\upharpoonright\eta_{1}$ for every
$\alpha\in X$ and $i<m.$

\item Player $\mathsf{II}$ has a winning strategy in the game \emph{G}%
$_{F}\left(  T,E\right)  .$
\end{enumerate}
\end{claim}

\qquad\ \ \ \qquad\ 

Indeed, the existence of $F$ and $\eta_{1}$ follows by Corollary
\ref{Corol del juego}. We can now find $\eta$ with the following properties:

\begin{enumerate}
\item $\eta_{1}<\eta<\delta.$

\item If $\alpha,\beta\in X$ and $t\in F,$ then \textsf{Tr}$(\triangle
(t,\rho_{0\alpha}),\beta)\cap\delta\subseteq\eta.$
\end{enumerate}

\qquad\ \qquad\ \ 

We now have the following:

\begin{claim}
There is $Y\in E\cap M$ with the following properties:
\label{Claim PFA C suc cosa Y}

\begin{enumerate}
\item $\eta<$ \textsf{min}$\left(  Y\right)  .$

\item Every element of $Y$ extends an element of $F.$

\item If $\gamma\in X$ and $\alpha,\beta\in Y$ are such that $\alpha<\beta,$
then \textsf{Tr}$(\alpha,\gamma)\cap M\subseteq\beta.$
\end{enumerate}
\end{claim}

\qquad\ \ \ \ \ 

Let $\sigma\in M$ be a winning strategy for Player $\mathsf{II}$ in the game
\emph{G}$_{F}\left(  T,E\right)  .$ Consider the following run of the game (in
which Player $\mathsf{II}$ is following her strategy\footnote{Note that in
order for $\sigma$ to be a winning strategy, it is necessary that all the
moves of Player $\mathsf{II}$ are of the form $\rho_{0\varepsilon}$ for
$\varepsilon<\omega_{1}$ (otherwise she inmediately loses the game).} $\sigma
$).\newline

\hfill%
\begin{tabular}
[c]{|l|l|l|l|l|l|}\hline
$\mathsf{I}$ & $\eta$ &  & $\xi_{0}$ &  & $...$\\\hline
$\mathsf{II}$ &  & $\rho_{0\varepsilon_{0}}$ &  & $\rho_{0\varepsilon_{1}}$ &
\\\hline
\end{tabular}
\hfill\ 

\qquad\ \qquad\ \ 

In which Player $\mathsf{I}$ plays as follows:

\begin{enumerate}
\item His first move is $\eta.$

\item After Player $\mathsf{II}$ played her move (say $\rho_{0\varepsilon_{0}%
}$), Player $\mathsf{I}$ chooses $\xi_{0}\in M$ that is larger than $\eta,$
$\varepsilon_{0}$ and $\bigcup\limits_{\gamma\in X}$\textsf{Tr}$(\varepsilon
_{0},\gamma)\cap M.$

\item After Player $\mathsf{II}$ played her move (say $\rho_{0\varepsilon_{1}%
}$), Player $\mathsf{I}$ chooses $\xi_{1}\in M$ that is larger than
$\varepsilon_{1}$ and $\bigcup\limits_{\gamma\in X}$\textsf{Tr}$(\varepsilon
_{1},\gamma)\cap M.$

\item[$\vdots$] $\vdots\hfill\vdots\hfill\vdots$
\end{enumerate}

\qquad\ \ \qquad

Since $\sigma$ is a winning strategy for Player $\mathsf{II,}$ it follows that
$Y=\left(  \varepsilon_{0},...,\varepsilon_{m-1}\right)  $ \ is in $E$ and
clearly has the desired properties. This finishes the proof of the claim.

\qquad\qquad\qquad

Fix $Y$ as in the claim. By elementarity., there exists $r\in D\cap M$ that is
a witnesses of $Y\in E.$ Let $h:S\left(  p\right)  \longrightarrow S\left(
r\right)  $ be the unique isomorphism from $H\left(  p\right)  $ to $H\left(
r\right)  .$ Recall that $h$ is the identity on $M.$ We claim that $p$ and $r$
are compatible. It suffices to show that $q=(\mathcal{M}_{p}\cup
\mathcal{M}_{r},\mathcal{N}_{p}\cup\mathcal{N}_{r})$ belongs to $\mathbb{P(}%
C\mathbb{)}.$ Clearly $\mathcal{N}_{p}\cup\mathcal{N}_{r}$ is a finite $\in
$-chain and contains $\mathcal{M}_{p}\cup\mathcal{M}_{r}.$

\qquad\ \qquad\ \ \ 

To verify that $q$ satisfies the completeness condition, it suffices to show
that for any $\alpha,\beta\in$ \textsf{ht}$_{\mathcal{M}}\left(  p\right)
\setminus M,$ we have that $\left[  h\left(  \alpha\right)  \beta\right]  \in
C.$ Denote $\triangle=\triangle\left(  h\left(  \alpha\right)  ,\beta\right)
.$ Recall that $h\left(  \alpha\right)  \in M$ while $\beta\notin M$, so
$h\left(  \alpha\right)  <\beta.$ It follows that $\left[  h\left(
\alpha\right)  \beta\right]  =$ \textsf{min}$($\textsf{Tr}$(\triangle
,\beta)\setminus h\left(  \alpha\right)  $). We proceed by cases, depending on
whether $\rho_{0\alpha}\upharpoonright\delta$ and $\rho_{0\beta}%
\upharpoonright\delta$ are equal.$\medskip$

$%
\begin{tabular}
[c]{|l|}\hline
\textbf{Case:} $\rho_{0\alpha}\upharpoonright\delta\neq\rho_{0\beta
}\upharpoonright\delta$ (equivalently, $\triangle\left(  \alpha,\beta\right)
<\delta$)\\\hline
\end{tabular}
\ \medskip$

In this case, we have that $\triangle=\triangle\left(  \alpha,\beta\right)  $
and \textsf{Tr}$(\triangle,\beta)\cap M\subseteq\eta_{0}\subseteq h\left(
\alpha\right)  .$ It follows that \textsf{Tr}$(\triangle,\beta)\setminus
h\left(  \alpha\right)  =$ \textsf{Tr}$(\triangle,\beta)\setminus\delta$ and
therefore \newline$\left[  h\left(  \alpha\right)  \beta\right]  =$
\textsf{min}$($\textsf{Tr}$(\triangle,\beta)\setminus\delta)$. The proof now
proceeds by subcases, depending on whether $\beta$ is equal to $\delta
$.\medskip

\textbf{Subcase: }$\rho_{0\alpha}\upharpoonright\delta\neq\rho_{0\beta
}\upharpoonright\delta$ and $\beta=\delta.$

\ \ \ \ \ 

We have that \textsf{Tr}$(\triangle,\beta)\setminus\delta=\left\{
\delta\right\}  ,$ so $\left[  h\left(  \alpha\right)  \beta\right]
=\delta\in C.\medskip$

\textbf{Subcase: }$\rho_{0\alpha}\upharpoonright\delta\neq\rho_{0\beta
}\upharpoonright\delta$ and $\beta\neq\delta.$

\qquad\qquad\qquad\ \ \ \ 

We know that $\left[  h\left(  \alpha\right)  \beta\right]  =$ \textsf{min}%
$($\textsf{Tr}$(\triangle,\beta)\setminus\delta),$ which is in $C$ by the
second crossing models condition of $p.\medskip$

$%
\begin{tabular}
[c]{|l|}\hline
\textbf{Case:} $\rho_{0\alpha}\upharpoonright\delta=\rho_{0\beta
}\upharpoonright\delta$ (equivalently, $\delta\leq\triangle\left(
\alpha,\beta\right)  $)\\\hline
\end{tabular}
\ \medskip$

The proof now proceeds by subcases, depending on the relationship of $\alpha$
and $\beta$ with $\delta$.\medskip

\textbf{Subcase:}\ \ \ $\rho_{0\alpha}\upharpoonright\delta=\rho_{0\beta
}\upharpoonright\delta$ and $\alpha=\delta$ or $\beta=\delta.$

\qquad\qquad\qquad\ \ 

Lemma \ref{Lema limite anticadena} implies that if $\rho_{0\alpha
}\upharpoonright\delta=\rho_{0\beta}\upharpoonright\delta$ and $\alpha=\delta$
or $\beta=\delta,$ then $\alpha=\beta=\delta.$ We have that \textsf{Tr}%
$(\triangle\left(  h\left(  \delta\right)  ,\delta\right)  ,\delta)\cap
\delta\subseteq\eta\subseteq h\left(  \delta\right)  ,$ so $\left[  h\left(
\delta\right)  \delta\right]  =\delta\in C.\medskip$

\textbf{Subcase:}\ \ \ $\rho_{0\alpha}\upharpoonright\delta=\rho_{0\beta
}\upharpoonright\delta$ and $\alpha,\beta\neq\delta.$

\qquad\ \qquad\ \qquad\ 

Recall that $\lambda\left(  \delta,\beta\right)  <\eta_{0}<\triangle.$
Proposition \ref{Prop juntar trazas} implies that $\delta\in$ \textsf{Tr}%
$(\triangle,\beta).$ Since \textsf{Tr}$(\triangle,\beta)\cap M\subseteq
\eta\subseteq h\left(  \alpha\right)  ,$ it follows that $\left[  h\left(
\alpha\right)  \beta\right]  =\delta\in C.$

\qquad\ \qquad\ \qquad\ \ \qquad\ \qquad\ \ 

This completes the proof that $q$ satisfies the completeness condition. We now
verify the first crossing models condition. It suffices to show that for every
$\alpha<\beta<\gamma$ with $\gamma\in$ \textsf{ht}$_{\mathcal{M}}\left(
p\right)  \setminus M$ and $\alpha,\beta\in$ \textsf{ht}$\left(  p\right)
\cup$ \textsf{ht}$\left(  r\right)  ,$ we have that \textsf{min}$($%
\textsf{Tr}$(\alpha,\gamma)\setminus\beta)\in C.$ The proof proceeds by
cases.$\medskip$

$%
\begin{tabular}
[c]{|l|}\hline
\textbf{Case:} \textbf{ }$\alpha\in$ \textsf{ht}$\left(  p\right)  $\\\hline
\end{tabular}
\ \medskip$

If $\beta$ is also in \textsf{ht}$\left(  p\right)  ,$ then there is nothing
to do since $p\in\mathbb{P}\left(  C\right)  .$ Assume that $\beta\in$
\textsf{ht}$\left(  r\right)  \setminus$ \textsf{ht}$\left(  p\right)  .$ Note
that \textsf{Tr}$(\alpha,\gamma)\cap M\subseteq\eta_{0}\subseteq\beta.$ It
follows that \newline\textsf{min}$($\textsf{Tr}$(\alpha,\gamma)\setminus
\beta)=$ \textsf{min}$($\textsf{Tr}$(\alpha,\gamma)\setminus\delta)$, which is
in $C$ by the first crossing models condition of $p.\medskip$

$%
\begin{tabular}
[c]{|l|}\hline
\textbf{Case:} $\alpha\in$ \textsf{ht}$\left(  r\right)  \setminus$
\textsf{ht}$\left(  p\right)  $ and $\beta\in$ \textsf{ht}$\left(  p\right)
$\\\hline
\end{tabular}
\ \medskip$

Note that $\delta\leq\beta.$ Since $\lambda\left(  \delta,\gamma\right)
<\eta_{0}<\alpha,$ Proposition \ref{Prop juntar trazas} implies that
\textsf{Tr}$(\delta,\gamma)$ is an initial segment of \textsf{Tr}%
$(\alpha,\gamma).$ It follows that \newline\textsf{min}$($\textsf{Tr}%
$(\alpha,\gamma)\setminus\beta)=$ \textsf{min}$($\textsf{Tr}$(\delta
,\gamma)\setminus\delta)$, which is in $C$ by the first crossing models
condition of $p.\medskip$

$%
\begin{tabular}
[c]{|l|}\hline
\textbf{Case:} $\alpha\in$ \textsf{ht}$\left(  r\right)  \setminus$
\textsf{ht}$\left(  p\right)  $ and $\beta\in$ \textsf{ht}$\left(  r\right)
$\\\hline
\end{tabular}
\ \medskip$

Since $\lambda\left(  \delta,\gamma\right)  <\eta_{0}<\alpha,$ Proposition
\ref{Prop juntar trazas} implies that $\delta\in$ \textsf{Tr}$(\alpha
,\gamma).$ Furthermore, since \textsf{Tr}$(\alpha,\gamma)\cap M\subseteq\beta$
(see Claim \ref{Claim PFA C suc cosa Y}), it follows that \newline%
\textsf{min}$($\textsf{Tr}$(\alpha,\gamma)\setminus\beta)=\delta\in C.$

\qquad\qquad\qquad\qquad

This completes the proof that $q$ satisfies the first crossing models
condition. We now verify the second crossing models condition. It suffices to
show that for every $N\in$ $\mathcal{N}_{p}\cup\mathcal{N}_{r}$ and
$\alpha,\beta\in$ \textsf{ht}$_{\mathcal{M}}\left(  p\right)  \cup$
\textsf{ht}$_{\mathcal{M}}\left(  r\right)  $ are such that $\triangle
=\triangle\left(  \alpha,\beta\right)  \in N$ and $\delta_{N}<\beta,$
$\alpha\neq\delta_{N},$ we have that \textsf{min}$($\textsf{Tr}$(\triangle
,\beta)\setminus N)\in C.$ As the reader probably guessed, the proof proceeds
by cases.$\medskip$

$%
\begin{tabular}
[c]{|l|}\hline
\textbf{Case:} $N\in\mathcal{N}_{p}\cap\mathcal{N}_{r}$ (equivalently,
$\mathcal{N\in N}_{p}\cap M$)\\\hline
\end{tabular}
\ \medskip$

If either $\alpha,\beta\in$ \textsf{ht}$\left(  p\right)  $ or $\alpha
,\beta\in$ \textsf{ht}$\left(  r\right)  ,$ then there is nothing to do since
both $p$ and $r$ are conditions, so we assume this is not the case.\medskip

\textbf{Subcase:} $N\in\mathcal{N}_{p}\cap\mathcal{N}_{r},$ $\alpha\in$
\textsf{ht}$\left(  p\right)  \setminus$ \textsf{ht}$\left(  r\right)  $ and
$\beta\in$ \textsf{ht}$\left(  r\right)  \setminus$ \textsf{ht}$\left(
p\right)  .$

\qquad\ \qquad\ \qquad\ \ 

Note that $\alpha\notin M.$ Moreover, since $\triangle\in N,$ it follows that
$\triangle<\eta_{0},$ which implies that $\triangle=\triangle\left(  h\left(
\alpha\right)  ,\beta\right)  .$ In this way, \textsf{min}$($\textsf{Tr}%
$(\triangle,\beta)\setminus N)=$ \textsf{min}$($\textsf{Tr}$(\triangle\left(
h\left(  \alpha\right)  ,\beta\right)  ,\beta)\setminus N),$ which is in $C$
by the second crossing models condition of $r.\medskip$

\textbf{Subcase:} $N\in\mathcal{N}_{p}\cap\mathcal{N}_{r},$ $\alpha\in$
\textsf{ht}$\left(  r\right)  \setminus$ \textsf{ht}$\left(  p\right)  $ and
$\beta\in$ \textsf{ht}$\left(  p\right)  \setminus$ \textsf{ht}$\left(
r\right)  .$

\qquad\qquad\ \ \qquad\ \ 

Note that $\beta\notin M.$ Since $\triangle<\eta_{0},$ it follows that
$\triangle=\triangle(h^{-1}\left(  \alpha\right)  ,\beta).$ In this way,
\newline\textsf{min}$($\textsf{Tr}$(\triangle,\beta)\setminus N)=$
\textsf{min}$($\textsf{Tr}$(\triangle\left(  h^{-1}\left(  \alpha\right)
,\beta\right)  ,\beta)\setminus N),$ which is in $C$ by the second crossing
models condition of $p.\medskip$

$%
\begin{tabular}
[c]{|l|}\hline
\textbf{Case:} $N\in\mathcal{N}_{r}\setminus\mathcal{N}_{p}$\\\hline
\end{tabular}
\ \medskip$

If both $\alpha,\beta$ belong to \textsf{ht}$\left(  r\right)  ,$ then there
is nothing to do, so we assume this is not the case.\medskip

\textbf{Subcase: }$N\in\mathcal{N}_{r}\setminus\mathcal{N}_{p}$, $\alpha\in$
\textsf{ht}$\left(  r\right)  $ and $\beta\notin$ \textsf{ht}$\left(
r\right)  .$

\qquad\ \ 

Note that $\beta\notin M.$ Let $\overline{\alpha}=h^{-1}\left(  \alpha\right)
.$ First consider the case where $\rho_{0\beta}\upharpoonright\delta\neq
\rho_{0\overline{\alpha}}\upharpoonright\delta.$ It follows that
$\triangle=\triangle(\overline{\alpha},\beta)$ and \textsf{Tr}$(\triangle
,\beta)\cap M\subseteq\eta_{0}\subseteq\eta.$ Moreover, \textsf{min}%
$($\textsf{Tr}$(\triangle,\beta)\setminus N)=$ \textsf{min}$($\textsf{Tr}%
$(\triangle\left(  \overline{\alpha},\beta\right)  ,\beta)\setminus M).$ If
$\beta=\delta,$ then this minimum is $\delta,$ which is in $C.$ If
$\delta<\beta,$ the conclusion follows by the second crossing models of $p.$
We now assume that\ $\rho_{0\beta}\upharpoonright\delta=\rho_{0\overline
{\alpha}}\upharpoonright\delta.$ We have that $\lambda\left(  \delta
,\beta\right)  <\eta_{0}<\triangle.$ Proposition \ref{Prop juntar trazas}
implies that $\delta\in$ \textsf{Tr}$(\triangle,\beta).$ Moreover,
\newline\textsf{Tr}$(\triangle,\beta)\cap M\subseteq\eta\subseteq N.$ It
follows that \textsf{min}$($\textsf{Tr}$(\triangle,\beta)\setminus
N)=\delta\in C.\medskip$

\textbf{Subcase: }$N\in\mathcal{N}_{r}\setminus\mathcal{N}_{p}$, $\alpha
\notin$ \textsf{ht}$\left(  r\right)  $ and $\beta$ $\in$ \textsf{ht}$\left(
r\right)  .$

\qquad\ \ \ \qquad\ \qquad\ \ 

Note that $\alpha\notin M.$ Let $\overline{\beta}=h^{-1}\left(  \beta\right)
.$ First consider the case where $\rho_{0\alpha}\upharpoonright\delta\neq
\rho_{0\overline{\beta}}\upharpoonright\delta.$ It follows that $\triangle
=\triangle(h\left(  \alpha\right)  ,\beta).$ Moreover, \textsf{min}%
$($\textsf{Tr}$(\triangle,\beta)\setminus N)=$ \textsf{min}$($\textsf{Tr}%
$(\triangle\left(  h\left(  \alpha\right)  ,\beta\right)  ,\beta)\setminus
N),$ which is in $C$ by the second crossing models condition of $r.$ We now
assume that\ $\rho_{0\alpha}\upharpoonright\delta=\rho_{0\overline{\beta}%
}\upharpoonright\delta.$ Since $\lambda(\delta,\overline{\beta})\in M$ and $h$
is an isomorphism that is the identity in $M,$ we have the following:\newline

\hfill%
\begin{tabular}
[c]{lll}%
$\lambda(\delta,\overline{\beta})$ & $=$ & $h(\lambda(\delta,\overline{\beta
}))$\\
& $=$ & $\lambda(h\left(  \delta\right)  ,h(\overline{\beta}))$\\
& $=$ & $\lambda(h\left(  \delta\right)  ,\beta)$%
\end{tabular}
\hfill\ 

\qquad\ \qquad\ \qquad\ \qquad\ \ \ \ 

Since $\lambda(\delta,\overline{\beta})<\eta_{0}<\triangle,$ it follows that
$\lambda(h\left(  \delta\right)  ,\beta)<\triangle.$ Proposition
\ref{Prop juntar trazas} implies that \textsf{Tr}$(h\left(  \delta\right)
,\beta)$ is an initial segment of \textsf{Tr}$(\triangle,\beta).$ Since
$h\left(  \delta\right)  \leq\delta_{N},$ we get that \textsf{min}%
$($\textsf{Tr}$(\triangle,\beta)\setminus N)=$ \textsf{min}$($\textsf{Tr}%
$(h\left(  \delta\right)  ,\beta)\setminus N),$ which is in $C$ by the first
crossing models condition of $p.\medskip$

\textbf{Subcase: }$N\in\mathcal{N}_{r}\setminus\mathcal{N}_{p}$ and
$\alpha,\beta\notin$ \textsf{ht}$\left(  r\right)  .$

\qquad\ \qquad\ \ 

Since $\triangle\in N,$ it follows that $\triangle\in M,$ which implies that
$\triangle<\eta_{0}.$ Moreover, \textsf{Tr}$(\triangle,\beta)\cap
M\subseteq\eta_{0}\in N.$ It follows that \textsf{min}$($\textsf{Tr}%
$(\triangle,\beta)\setminus N)=$ \textsf{min}$($\textsf{Tr}$(\triangle
,\beta)\setminus M),$ which is in $C$ by the second crossing models condition
of $p.\medskip$

$%
\begin{tabular}
[c]{|l|}\hline
\textbf{Case:} $N\in\mathcal{N}_{p}\setminus\mathcal{N}_{r}$\\\hline
\end{tabular}
\ \medskip$

Since $\delta_{N}<\beta,$ it follows that $\beta\in$ \textsf{ht}$\left(
p\right)  \setminus$ \textsf{ht}$\left(  r\right)  .$ If $\alpha$ is also in
\textsf{ht}$\left(  p\right)  ,$ then there is nothing to do, so assume that
$\alpha\in$ \textsf{ht}$\left(  r\right)  \setminus$ \textsf{ht}$\left(
p\right)  .$ Let $\overline{\alpha}=h^{-1}\left(  \alpha\right)  .$ If
$\rho_{0\overline{\alpha}}\upharpoonright\delta\neq\rho_{0\beta}%
\upharpoonright\delta,$ then $\triangle=\triangle(\overline{\alpha},\beta),$
so \textsf{min}$($\textsf{Tr}$(\triangle,\beta)\setminus N)$ is equal to
\textsf{min}$($\textsf{Tr}$(\triangle(\overline{\alpha},\beta),\beta)\setminus
N)$ and it is in $C$ by the second crossing models condition of $p.$ Assume
now that $\rho_{0\overline{\alpha}}\upharpoonright\delta=\rho_{0\beta
}\upharpoonright\delta.$ It follows that $\lambda(\delta_{N},\beta
)<\triangle.$ Proposition \ref{Prop juntar trazas} implies that $\delta_{N}%
\in$ \textsf{Tr}$(\triangle,\beta),$ so \textsf{min}$($\textsf{Tr}%
$(\triangle,\beta)\setminus N)=\delta_{N}\in C.$ At last we finished the proof.
\end{proof}

\end{document}
